\documentclass[sigconf,anonymous,review]{acmart}

%% ---------------------------------------------------------------
%% Compile with pdfLaTeX (Overleaf default). acmart then loads the
%% official ACM fonts (Linux Libertine + newtxmath).
\makeatletter
\@ifundefined{mathbb}{\usepackage{amsfonts}}{}
\makeatother
%% ---------------------------------------------------------------

%% Allow slightly looser lines instead of overfull boxes.
\setlength{\emergencystretch}{1.5em}
\usepackage{booktabs}
\usepackage{multirow}
\usepackage{makecell}
\usepackage{graphicx}

\setcopyright{none}
\copyrightyear{2027}
\acmYear{2027}
\acmConference[WWW '27]{The ACM Web Conference 2027}{May 10--14, 2027}{Dublin, Ireland}
\acmBooktitle{Proceedings of The ACM Web Conference 2027 (WWW '27)}
\acmDOI{}
\acmISBN{}
\settopmatter{printacmref=true}


%% ---------------------------------------------------------------
%% Review marks (yellow background, red bold).
%%   \tbd{...}  placeholder that MUST be filled before submission;
%%              always highlighted.
%%   \ph        numeric placeholder (prints X.XX); MUST be replaced.
%%   \ai{...}   wording flagged for human review; set \reviewmarksfalse
%%              to print it as normal text once reviewed.
%% Before submission: searching the project for "\tbd{" and "\ph"
%% must return only these definitions.
\usepackage{xcolor}
\newif\ifreviewmarks
\reviewmarkstrue
\ExplSyntaxOn
\seq_new:N \l__rev_words_seq
\bool_new:N \l__rev_first_bool
\NewDocumentCommand{\revhl}{m}
  {
    \group_begin:
    \setlength{\fboxsep}{0.6pt}
    \color{red}\bfseries
    \seq_set_split:Nnn \l__rev_words_seq {~} {#1}
    \bool_set_true:N \l__rev_first_bool
    \seq_map_inline:Nn \l__rev_words_seq
      {
        \bool_if:NTF \l__rev_first_bool
          { \bool_set_false:N \l__rev_first_bool }
          { \colorbox{yellow}{\strut\hspace{0.25em}}\hskip 0pt plus 0.15em\relax }
        \colorbox{yellow}{\strut ##1}
      }
    \group_end:
  }
\ExplSyntaxOff
\newcommand{\tbd}[1]{\revhl{#1}}
\newcommand{\ph}{\revhl{X.XX}}
\newcommand{\ai}[1]{\ifreviewmarks\revhl{#1}\else#1\fi}
%% ---------------------------------------------------------------

\newcommand{\method}{Omni-Aware}
\newcommand{\Uomni}{U_{\mathrm{omni}}}
\newcommand{\hgrd}{\mathcal{H}_{\mathrm{grd}}}
\newcommand{\hpred}{\mathcal{H}_{\mathrm{pred}}}
\newcommand{\taub}{\tau_{\mathrm{bias}}}
\newcommand{\Mbias}{M_{\mathrm{bias}}}
\newcommand{\tans}{t_{\mathrm{ans}}}
\newcommand{\tend}{t_{\mathrm{end}}}
\newcommand{\webimg}{WebQA-img}

\begin{document}


\title[Omni-Aware: Latent-Space Evidence Control]{Omni-Aware: Latent-Space Evidence Control for Multimodal Web Answering}

\author{Anonymous Author(s)}
\affiliation{\institution{Anonymous Institution}\country{}}

\begin{abstract}
Web retrieval-augmented generation answers questions from evidence that may lie in text passages or in images.
For each question a system makes three decisions: whether to retrieve, which modality to retrieve from, and how the retrieved evidence is used.
Existing multimodal systems make them with separate, format-specific signals, trained routers, or sampled generations.
We study \method{}, a controller that makes all three from one family of latent-space signals: the first-token distributions that a frozen generator produces when it prefills a prompt without evidence and with each kind of evidence.
We first examine the soft form of such control, Evidence Attention Bias, which adds a relevance-scaled bias to the attention over evidence spans.
On the previous pipeline it moved attention onto supporting evidence and left the answers almost unchanged, because the first-token margin exceeded the logit displacement that the bias produced.
\method{} therefore controls evidence in hard form: it decides which evidence set the generator reads.
Its rule, v2-Rule, fits no weights. It keeps the direct answer when the direct first-token margin clears a threshold and otherwise hands the generator text, image, mixed, or reranked evidence; applying it needs prefills and one decode.
A calibrated variant, v2-Cal, fits one small ridge regression per evidence action.
Thresholds are screened by fixed-sequence one-sided paired tests on the mean score gap to the strongest single action, with $\delta=0.10$ and a tolerance $\epsilon=1.0$ point: with probability at least $1-\delta$ over the calibration draw, the expected score of the chosen policy is at least that of the strongest single action minus $\epsilon$.
The strongest single action is chosen on the same calibration sample, which is slightly optimistic.
All actions share one frozen answer prompt; OK-VQA is scored with official VQA accuracy, and WebQA primarily with acc$_{\mathrm{approx}}$, with normalized BART and QA$_{\mathrm{approx}}$ on image and text subsets.
The evaluation of the repaired pipeline has not been run; every new experimental number in this draft is a placeholder.
\end{abstract}

\begin{CCSXML}
<ccs2012>
<concept><concept_id>10002951.10003317.10003338</concept_id><concept_desc>Information systems~Retrieval models and ranking</concept_desc><concept_significance>500</concept_significance></concept>
<concept><concept_id>10010147.10010178.10010179.10010182</concept_id><concept_desc>Computing methodologies~Natural language generation</concept_desc><concept_significance>300</concept_significance></concept>
</ccs2012>
\end{CCSXML}
\ccsdesc[500]{Information systems~Retrieval models and ranking}
\ccsdesc[300]{Computing methodologies~Natural language generation}

\keywords{Web retrieval-augmented generation, multimodal question answering, latent-space control, adaptive retrieval, evidence use, risk control}

\maketitle

%% ===============================================================
\section{Introduction}
\label{sec:intro}

Web retrieval-augmented generation answers a question by retrieving evidence and conditioning a generator on it~\cite{guu2020realm,lewis2020rag}.
On the Web the supporting facts may sit in a text passage, an image, or both, so a multimodal system makes three decisions for every question: whether to retrieve at all, which modality to retrieve from, and how the retrieved evidence is used by the generator.
Current systems make these decisions with different machinery.
Whether to retrieve is decided from generated text, sampled generations, or a trained classifier; where to retrieve is decided by a router trained on the query; and how evidence is used is left to the generator's attention, which is sensitive to position and to the amount of evidence.
The three are therefore hard to coordinate, and each new format brings its own rule.

We take a latent-space view of the same decisions.
Evidence reaches a decoder-only generator as token spans in its input, and one forward pass over a prompt, a prefill with no generation, already yields the distribution of the first answer token.
Comparing that distribution without evidence and with each kind of evidence shows how confident the generator is on its own and what each evidence set would change.
\method{} makes all three decisions from this one family of signals, on a frozen generator, with the retrievers and encoders of each format left as they are.

The most direct way to act on such signals is soft control: bias the attention toward the evidence that the retriever scores highest.
We built this as Evidence Attention Bias and measured it on the previous pipeline.
The bias did move attention onto supporting evidence, and the answers barely changed: the margin of the first answer token was larger than the logit displacement that the bias produced, and among the evidence already selected the relevance scores were too close for the bias to separate them.
Section~\ref{sec:steer} reports that analysis, which comes from the project root-cause report dated 13~October~2026 and is not a result of the new runs.
Reweighting evidence that the generator already reads does not change what it concludes.

\method{} therefore controls evidence in hard form: it decides which evidence set the generator reads.
The action set contains a direct answer and four evidence sets, text only, image only, mixed, and reranked, and it spans the three decisions: direct against the rest is whether to retrieve, text against image is where, and mixed against reranked is how the evidence is selected.
A rule with no fitted weights keeps the direct answer when its first-token margin clears a threshold and otherwise selects an evidence action; it is applied with prefills and a single decode.
Its thresholds are screened by a paired test on the calibration split, which bounds the expected loss against the strongest single action (Section~\ref{sec:risk}).

Format and metric choices can move scores by more than the gap between systems.
The evaluation therefore freezes one answer prompt for every action, scores OK-VQA with official VQA accuracy only, and scores WebQA primarily with acc$_{\mathrm{approx}}$, reporting normalized BART and QA$_{\mathrm{approx}}$ on image and text subsets beside it.
The primary relative gain uses, as its denominator, the single action with the best score on the calibration split.
Historical cached baselines are not used.

Our contributions are as follows:
\begin{itemize}
  \item A latent-space formulation of evidence control for multimodal Web answering, in which whether to retrieve, where to retrieve, and how evidence is used are decided from the first-token distributions of prefills with and without evidence, on a frozen generator.
  \item An analysis of soft control (Section~\ref{sec:steer}): a relevance-scaled attention bias moves attention onto supporting evidence and does not move answers, with the first-token margin as the explanation.
  Every number in that subsection is copied from the 13~October~2026 root-cause report and is attributed there.
  \item Hard control: v2-Rule, which fits no weights and is applied with prefills and one decode, and v2-Cal, one ridge regression per evidence action.
  Both are screened by fixed-sequence one-sided paired tests on the mean score gap to the strongest single action (Section~\ref{sec:risk}), at $\delta=0.10$ and $\epsilon=1.0$; the guarantee concerns expected score, under the assumptions stated there.
  \item A measurement, under one frozen answer format and official metrics, of when retrieval helps or hurts on WebQA and OK-VQA, including image and text subsets and the gap between the best single action and the per-question oracle.
  The new measurement has not been run.
  Its cells, including that gap, are \ph{}.
\end{itemize}

%% ===============================================================
\section{Related Work}
\label{sec:related}

\textit{Web and multimodal retrieval-augmented generation.}
Web answer engines retrieve and cite pages at query time~\cite{nakano2021webgpt,liu2023webglm}, and WebQA poses multimodal Web answering over text and image facts~\cite{chang2022webqa}.
Multimodal RAG systems~\cite{abootorabi2025survey} retrieve images, knowledge memories, or visual documents~\cite{adjali2024multilevel,chen2022murag,hu2023reveal,wang2025vidorag,yu2025visrag}, generate multimodal answers~\cite{ma2025m2rag,wu2025visualrag}, or refine selection by reranking, reflection, or adaptive visual selection~\cite{chen2025ragvl,zhai2025samrag,zhang2024mr2ag}.
They improve which chunks are retrieved and how each format is encoded.
We keep those format-specific retrievers and encoders.
The controller in this paper decides which of their evidence sets the generator reads (Section~\ref{sec:problem}).

\textit{Adaptive evidence acquisition.}
Adaptive retrieval decides from the generation state whether external knowledge is needed: Self-RAG learns reflection signals~\cite{asai2024selfrag}, FLARE retrieves when upcoming tokens have low probability~\cite{jiang2023flare}, DRAGIN combines token entropy with attention to decide when and what to retrieve~\cite{su2024dragin}, Adaptive-RAG classifies question complexity~\cite{jeong2024adaptiverag}, and SeaKR measures the consistency of sampled hidden states~\cite{yao2025seakr}; their cost--quality tradeoffs have been compared~\cite{moskvoretskii2025adaptive}.
For multimodal inputs, MMA-RAG classifies hidden states of the input and the retrieved visual context~\cite{du2026mmarag}, and UniversalRAG routes the query alone to modality- and granularity-specific corpora~\cite{yeo2026universalrag}; retriever signals do not reliably improve such query-only routers~\cite{li2026beyond}.
These methods read generated text, the query, retriever scores, or pooled hidden states, and most train a router or classifier or sample many generations.
\method{} reads the first-token distributions of prefills with and without each evidence set: the signal is taken inside the generator, before any text is produced, and the same signal serves the whether, where, and how decisions. Its rule fits no weights and samples no generations (Section~\ref{sec:rule}), and its thresholds are screened with the procedure in Section~\ref{sec:risk}.

\textit{Knowledge conflict and inference-time evidence control.}
Retrieved context can agree or conflict with parametric knowledge, which is why the action set contains a direct arm and separate text, image, and mixed arms.
Position sensitivity~\cite{liu2024lost} has motivated multi-scale positional encoding~\cite{zhang2024mspoe} and positional attention calibration~\cite{hsieh2024found}; text-only RAG readers are less affected than long-context models~\cite{cuconasu2025positional}, and multimodal readers can prefer early passages, while reranking does not by itself remove that preference~\cite{liu2026lostend}.
OpenDecoder multiplies attention scores by retrieval features and fine-tunes the decoder~\cite{mo2026opendecoder}; PASTA emphasizes user-specified spans on profiled heads~\cite{zhang2024pasta}, which at one attention row equals a constant span bias; PAI amplifies attention to image tokens~\cite{liu2024pai}; and MAD-RAG restores attention to question-relevant image regions that retrieved text draws away~\cite{zhao2026madrag}.
Section~\ref{sec:steer} quotes a measurement in which an attention bias moves attention onto supporting evidence while the answer stays put.
These methods are soft control in our terms: they reweight evidence that the generator already reads. Evidence Attention Bias belongs to the same family, and that analysis is why \method{} selects the evidence set instead.

%% ===============================================================
\section{Problem Definition}
\label{sec:problem}

An instance is a query $q$, optional context such as the OK-VQA question image, and an evidence pool.
A frozen generator produces an answer string.
The main prompt is shared by every action.
It asks for the frozen answer format and does not instruct the model to emit a sources line in place of an answer.
WebQA answers are one sentence aligned with the reference sentence style.
OK-VQA answers are at most three words after official VQA normalization.
That format is fixed before any model run and is the same for every arm.

The action set is
\begin{equation}
\begin{aligned}
\mathcal{A}=\{\texttt{direct},\ &\texttt{static\_text},\ \texttt{static\_image},\\
&\texttt{static\_mixed},\ \texttt{rerank\_static}\}.
\end{aligned}
\label{eq:actions}
\end{equation}
\texttt{direct} uses no retrieved evidence.
The three static actions pack text, image, or mixed evidence.
\texttt{rerank\_static} packs the top-$k$ of the whole candidate pool, ranked by a cross-encoder reranker.
The set spans the three decisions: \texttt{direct} against the others is whether to retrieve; \texttt{static\_text} against \texttt{static\_image} is where; and \texttt{static\_mixed} against \texttt{rerank\_static} is how the evidence is selected.
Evidence Attention Bias is soft control and is not a member of $\mathcal{A}$.
All five actions are decoded in the same run, on the same candidates, with the same prompt and the same decoding budget.

The selector returns one action $a\in\mathcal{A}$.
Let $s(\cdot)$ be the official metric on a given split, and let $a^{\star}$ be the single action in $\mathcal{A}$ with the highest score on the calibration split.
The primary relative gain on a later split uses $s(a^{\star})$ as the denominator.
The same tables also report the gain versus \texttt{static\_mixed}.
The oracle picks, for each question, the action in $\mathcal{A}$ with the best score.
Risk control is stated in Section~\ref{sec:risk}.
Its claim is a bound on the expected score of the chosen policy relative to $a^{\star}$.

%% ===============================================================
\section{Method}
\label{sec:method}

\subsection{Soft and Hard Evidence Control}
\label{sec:softhard}
Evidence enters the generator as token spans, and a prefill over a prompt gives the distribution of the first answer token without generating anything.
A controller can act on that information in two ways.
Soft control leaves the input as it is and changes how strongly the generator attends to each evidence span; Evidence Attention Bias adds $\gamma U r_j$ to the attention logits of the columns of evidence chunk $j$, with $r_j$ its calibrated relevance and $U$ an uncertainty score.
Hard control changes the input: it chooses which evidence set is packed into the prompt.
Section~\ref{sec:steer} shows why soft control is not the main method here.
On the previous pipeline the bias shifted attention mass toward supporting chunks, and the first answer token rarely changed, because its margin over the runner-up was larger than the displacement the bias caused, and because the chunks already selected had relevance scores too close to be told apart.
The rest of this section is hard control over the action set $\mathcal{A}$.

\subsection{Signals}
\label{sec:signals}
Each action stores the entropy of its first-token distribution and the same entropy divided by $\log|V|$, the gap between the top-1 and top-2 logits, the log-probability of the greedy answer and the number of generated tokens, and, for an evidence action, the KL divergence between that action's first-token distribution and the direct action's first-token distribution.
The run also stores the retriever's top-1 score and the gap between retriever top-1 and top-2, per modality.
\texttt{direct} has no KL term; its retriever fields describe the pool's top-1 score, not a selected list.
\texttt{gold-in-top-$k$} is computed only for offline analysis.
It is not an inference-time feature and it is not an input to v2-Cal.

Only four quantities enter a decision: the first-token margin, the token F1 between an action's answer and the answer of $a^{\star}$, the KL divergence from \texttt{direct}, and the retriever top-1 score of the candidates an action uses. Entropy, the sequence log-probability, the generated length, and the retriever top-1--top-2 gap are stored for analysis.

The first-token signals of an action need only a prefill of its prompt.
The decoded answer, and therefore the token F1, needs a decode.
v2-Rule reads first-token margins only, so it can be applied before anything is generated (Section~\ref{sec:cost}); v2-Cal also reads token F1 and needs the decoded answers.

\subsection{v2-Rule}
\label{sec:rule}
v2-Rule fits no weights.
Let $c(\cdot)$ be the first-token margin, the gap between the top-1 and top-2 logits.
If $c(\texttt{direct})\ge\tau$, the rule selects \texttt{direct}.
Otherwise it selects an evidence action in one of two ways: the evidence action $e^{\star}$ with the highest calibration mean, which equals $a^{\star}$ whenever $a^{\star}$ is an evidence action, or the available evidence action with the largest margin on that question.
Evidence actions with an empty candidate pool are unavailable and are skipped.
The thresholds $\tau$ are the 0\%, 5\%, \dots, 100\% quantiles of the direct margin on the calibration split, together with $\pm\infty$.
The sequence starts at $\tau=-\infty$ when $a^{\star}$ is \texttt{direct} and at $\tau=+\infty$ when it is an evidence action, and moves away from there. With the best-evidence rule this starting policy equals $a^{\star}$; with the most-confident rule it equals $a^{\star}$ only when $a^{\star}$ is \texttt{direct}, and otherwise the whole sequence is rejected if its first policy fails the test.
Each evidence rule has its own sequence at level $\delta/2$.
The threshold is chosen by utility (Section~\ref{sec:risk}), not by accuracy.

\subsection{v2-Cal}
\label{sec:cal}
v2-Cal fits one ridge regression for each evidence action $b\neq a^{\star}$.
Its target is the score gap $s(b)-s(a^{\star})$ on the fitting half of the calibration split.
Each regression has an intercept and six features: the direct margin, the margin of $a^{\star}$, the margin gap $c(b)-c(a^{\star})$, the token F1 between the answers of $b$ and $a^{\star}$, the KL divergence of $b$ from \texttt{direct}, and the retriever top-1 score of $b$.
The intercept is not penalized, and the standardized features receive a ridge penalty of 1.0; the coefficients are returned in raw units.
A feature missing on every fitting question, or constant on all of them, is dropped and recorded in the model.
A feature missing on only some questions is an error.
Missing values are never filled with zero.
The policy takes the evidence action with the largest predicted gain when that gain exceeds a threshold $\lambda_{\mathrm{gain}}$, and $a^{\star}$ otherwise.
The calibration split is halved with a fixed seed: the fitting half trains the regressions, and the other half runs the risk control of Section~\ref{sec:risk}.
Unavailable actions are excluded from fitting and from comparison.
We describe v2-Cal as calibration, not as a trained router.
A previous OK-VQA rule (\texttt{original\_retrieve} or $p_{\mathrm{image}}<0.5$) used a trained query classifier.
v2 does not use that classifier.

\subsection{Risk Control}
\label{sec:risk}
Both selectors are screened by a fixed-sequence test on the mean score gap to $a^{\star}$.
For a calibration question $i$ and a policy $\pi$, let $d_i=s_i(\pi)-s_i(a^{\star})$ be the paired gap.
The null hypothesis is $H_0:\mathbb{E}[d]<-\epsilon$, tested with the one-sided paired statistic
\begin{equation}
t=\frac{\bar d+\epsilon}{\hat\sigma_d/\sqrt{n}},
\label{eq:stat}
\end{equation}
whose $p$-value is taken from the standard normal upper tail at every sample size.
A constant gap has zero variance: its $p$-value is 0 when its mean is at least $-\epsilon$, and 1 otherwise.
The candidates are walked in their fixed order, and a candidate is accepted while its $p$-value is at most its family's level: $\delta/2$ for each evidence rule, and $\delta$ for v2-Cal.
Among the accepted candidates, the one with the largest mean gap $\bar d$ is chosen.
If none is accepted, or the largest mean gap is not positive, the selector returns $a^{\star}$.
The two parameters have separate roles: $\delta=0.10$ is the failure probability, and $\epsilon=1.0$ point is the tolerated loss in expected score.
Under i.i.d. questions and an asymptotically valid test, with probability at least $1-\delta$ over the calibration draw, the chosen policy satisfies
\begin{equation}
\mathbb{E}[s(\pi)]\ \ge\ \mathbb{E}[s(a^{\star})]-\epsilon.
\label{eq:guarantee}
\end{equation}
The guarantee is a statement about expected score over questions, not a per-question or per-split statement, and it does not claim that the method is never worse.
The strongest single action is chosen on the same calibration sample, which makes the guarantee slightly anti-conservative.
v2-Cal is tested at level $\delta$ as its own decision, and the main text reports v2-Rule as the primary one.

\subsection{Cost}
\label{sec:cost}
Calibration and evaluation decode every action of every question, because the score of each action is needed to choose thresholds and to report the oracle.
Applying the frozen v2-Rule needs less.
It prefills the direct prompt, one forward pass with no generation.
If the direct margin clears $\tau$, it decodes the direct answer.
Otherwise the best-evidence rule decodes $e^{\star}$ with no further prefill, and the most-confident rule first prefills each available evidence action.
A question therefore costs at most $|\mathcal{A}|$ prefills and exactly one decode.
v2-Cal reads decoded answers and has no such form.
The test runs execute this path beside the full decode and record, per question, the prefills, the decode, the wall time, and whether the path chose the action that the full decode chooses.
The evaluation reports those counts, the number of input tokens, and the latency, next to always decoding \texttt{static\_mixed} and next to SeaKR~\cite{yao2025seakr}.
Those measurements are \ph{}.
The GPU-time figures in the repository README are planning estimates, not measurements, and they are not copied into the tables.

%% ===============================================================
\section{Experiments}
\label{sec:experiments}
The v2 evaluation has not been run.
Every new experimental number below is \ph{}.
Numbers copied from the root-cause report dated 13~October~2026 are marked as such.
They were computed on the previous pipeline and are not v2 results.
Exploratory outputs of \texttt{offline\_v2} are from that old pipeline, contain known bugs, and are not the main table.
The section is a measurement study, a main table, ablations, and analysis.
No cell has been filled from a v2 run.

\subsection{Experimental Setup}
\label{sec:setup}
\textit{Splits.}
WebQA calibration is at least 500 questions and the development split is 256 questions.
Both are drawn from the training split, they are disjoint, and neither is drawn from the validation split.
The test split is all 4{,}966 validation questions.
The previous statement that calibration came from the validation split, while the reported set was the entire validation split, is withdrawn.
OK-VQA uses 500 calibration questions, 256 development questions, and the remainder of the validation split as the test set.
The three OK-VQA splits are pairwise disjoint.
An empty test split is an error.
When the validation split has 5{,}046 questions, the remainder is 4{,}290.
MMQA and video are generalization settings whose sizes follow the available assets.
The current command interface does not include them, and their scores are \ph{}.

\textit{Settings and metrics.}
The primary WebQA setting is the distractor pool (\texttt{webqa\_full}), in which each question brings its own candidates.
Open retrieval (\texttt{webqa\_open\_*}; the current interface generates \texttt{webqa\_open\_img} at stage~s2) is the second setting.
OK-VQA is scored with official VQA accuracy only.
WebQA's primary metric is acc$_{\mathrm{approx}}$.
Normalized BART and QA$_{\mathrm{approx}}$ are reported as well, on the full set and on the image and text subsets.
The image subset has 2{,}511 questions and the text subset has 2{,}455.
Loose lexical WebQA scores are not placed next to acc$_{\mathrm{approx}}$.

\textit{Actions and the strong baseline.}
The decoded actions are \texttt{direct}, \texttt{static\_text}, \texttt{static\_image}, \texttt{static\_mixed}, and \texttt{rerank\_static}.
Adaptive comparators are an NTE trigger, SeaKR, the previous $\Uomni$ gate (v1), and UniversalRAG routing.
They are not eligible to be the strong baseline.
The strong baseline is the single action with the best calibration-set score.
Primary relative gain uses that score as the denominator.
The tables also report gain versus \texttt{static\_mixed}.
The oracle is the per-question best action in the same set.
Evidence Attention Bias is not a main-method arm.
All actions in a run share one prompt, one budget, and one candidate pool.
The prompt is fixed before any model run.
The test split is scored once, after the frozen decision file \texttt{runs/v2\_decisions/} is written.
Grids are not selected on the test split.
The operating point previously taken from a test-set grid (Table~20 in the audited draft) is withdrawn.

\textit{Statistics.}
Uncertainty is a paired bootstrap 95\% confidence interval over questions.
Multiple comparisons use Holm correction.
Decoding is greedy, so a configuration is deterministic given its inputs.
This draft does not report a seed standard deviation under greedy decoding.
The risk-control claim is reported as the mean paired gap to the strong baseline, with its bootstrap interval.
A score below \texttt{direct} is not described as better than no retrieval.

\textit{Development gates.}
Two gates, written down before any GPU run of the repaired pipeline, are criteria, not results, and this draft does not claim that either has passed.
After the pipeline fixes, the development-set per-question oracle is compared with the strongest single action against a relative gap of 8\%.
A selector is eligible for the frozen test run when, on the development split, its relative gain over that baseline is at least 5\%, and the lower end of its paired interval is positive.
The test split is not used to make those choices.

\textit{Figures withdrawn by the root-cause report.}
The following numbers are copied from the root-cause report of 13~October~2026.
They come from the legacy pipeline, which has known prompt and format problems; they motivate the protocol and are not results of the method. They are not v2 results, and they are not written into Tables~\ref{tab:webqa}--\ref{tab:cost}.
Relative gains of +33\% on WebQA (official QA score 18.1 to 24.2 on the 4{,}966 validation questions) and +15\% on OK-VQA compared against historical cached baselines produced under different prompts, budgets, and retrieval scores. On one 512-question batch the cached static score was 17.9 where static retrieval decoded in the same run scored 23.4, a gap of 5.5 points from the protocol alone.
The same report withdraws a trained controller's +8.43\% and a +14.4\% chosen on a 32-question development set.
On a 256-question OK-VQA development slice, official VQA accuracy was 61.5 for direct answering, 36.3 for static retrieval, 48.1 for the gated static arm, and 36.3 for EAB.
Direct answering won 82 questions and lost 10 against static retrieval.
The per-question maximum of direct and static on that slice was 64.4, which is 2.9 points above direct answering (about +4.7\% relative).
The gated arm and EAB are both below direct answering, so the report does not call that slice an improvement over no retrieval.
On 512 WebQA questions, static retrieval scored 59.1 acc$_{\mathrm{approx}}$, 27.4 normalized BART, and 23.4 QA.
Image questions scored 63.2 acc$_{\mathrm{approx}}$ and 13.8 QA.
Text questions scored 55.0 acc$_{\mathrm{approx}}$ and 33.0 QA.
Over 2{,}560 items, a loose WebQA scorer and official acc$_{\mathrm{approx}}$ differed by 10.5 points on average (0.699 versus 0.594).
On one static OK-VQA batch, a loose scorer gave 65.8 where official VQA accuracy was about 47.9.
A greedy-decoding figure of $\pm 0.2$ is smaller than the sampling standard error of about 0.66 points that 5{,}046 OK-VQA questions imply, and it is not a confidence interval.
The v2 measurement that replaces these comparisons is the placeholder tables below.

\subsection{Measurement Study}
\label{sec:measure}
The measurement reserves three quantities.
The gap between the best single action and the per-question oracle is \ph{}.
Image questions and text questions are separate columns in Tables~\ref{tab:webqa} and~\ref{tab:okvqa}, and every column is \ph{}.
The size of the normalized BART factor, reported beside acc$_{\mathrm{approx}}$, is \ph{}.
A 50-question image-evidence diagnostic is specified, and its conclusion is \ph{}.
The audit's 256-question and 512-question figures above are the motive for these tables.
They are not substitutes for the cells.

\begin{table*}[t]
\caption{WebQA distractor test split (\texttt{webqa\_full}): placeholder measurement.
acc$_{\mathrm{approx}}$ is primary.
Normalized BART and QA$_{\mathrm{approx}}$ are secondary.
Image and text are subsets of the same run.
The strong baseline is whichever single action scores highest on the calibration split; its identity is frozen before test and is not chosen on this table.
The oracle is the per-question best action.
Every entry is \ph{}.
The v2 run has not been executed.}
\label{tab:webqa}
\scriptsize
\setlength{\tabcolsep}{3.5pt}
\begin{tabular}{@{}lccccccccc@{}}
\toprule
Action & \multicolumn{3}{c}{acc$_{\mathrm{approx}}$} & \multicolumn{3}{c}{norm.\ BART} & \multicolumn{3}{c}{QA$_{\mathrm{approx}}$} \\
\cmidrule(lr){2-4}\cmidrule(lr){5-7}\cmidrule(l){8-10}
 & All & Image & Text & All & Image & Text & All & Image & Text \\
\midrule
\texttt{direct} & \ph{} & \ph{} & \ph{} & \ph{} & \ph{} & \ph{} & \ph{} & \ph{} & \ph{} \\
\texttt{static\_text} & \ph{} & \ph{} & \ph{} & \ph{} & \ph{} & \ph{} & \ph{} & \ph{} & \ph{} \\
\texttt{static\_image} & \ph{} & \ph{} & \ph{} & \ph{} & \ph{} & \ph{} & \ph{} & \ph{} & \ph{} \\
\texttt{static\_mixed} & \ph{} & \ph{} & \ph{} & \ph{} & \ph{} & \ph{} & \ph{} & \ph{} & \ph{} \\
\texttt{rerank\_static} & \ph{} & \ph{} & \ph{} & \ph{} & \ph{} & \ph{} & \ph{} & \ph{} & \ph{} \\
Strong baseline & \ph{} & \ph{} & \ph{} & \ph{} & \ph{} & \ph{} & \ph{} & \ph{} & \ph{} \\
Oracle & \ph{} & \ph{} & \ph{} & \ph{} & \ph{} & \ph{} & \ph{} & \ph{} & \ph{} \\
\bottomrule
\end{tabular}
\end{table*}

\begin{table}[t]
\caption{OK-VQA test split: official VQA accuracy only.
Placeholder measurement.
The same action definitions as Table~\ref{tab:webqa}.
Every entry is \ph{}.}
\label{tab:okvqa}
\small
\begin{tabular}{@{}lc@{}}
\toprule
Action & VQA accuracy \\
\midrule
\texttt{direct} & \ph{} \\
\texttt{static\_text} & \ph{} \\
\texttt{static\_image} & \ph{} \\
\texttt{static\_mixed} & \ph{} \\
\texttt{rerank\_static} & \ph{} \\
Strong baseline & \ph{} \\
Oracle & \ph{} \\
\bottomrule
\end{tabular}
\end{table}

\subsection{Main Results}
\label{sec:main}
Table~\ref{tab:main} is the main comparison.
Relative gain versus the strong baseline is the primary column.
Gain versus \texttt{static\_mixed} is secondary.
Each gain cell is a point estimate and a paired bootstrap 95\% interval; both are \ph{}, so the cell is written \ph{} [\ph{},\,\ph{}].
The gap column is the mean paired score difference to the strong baseline, with its interval.
No row is marked significant, because no test has been run.
A filled row that sits below \texttt{direct} will not be described as better than no retrieval.

\begin{table*}[t]
\caption{Main results on the frozen test splits.
WebQA uses acc$_{\mathrm{approx}}$ in the distractor setting.
OK-VQA uses official VQA accuracy.
``Versus strong'' uses the calibration-set best single action as the denominator.
``Versus static'' uses \texttt{static\_mixed}.
Brackets are paired bootstrap 95\% intervals.
Gap is the mean paired score difference to the strong baseline.
All entries are \ph{}.
EAB is omitted.}
\label{tab:main}
\scriptsize
\setlength{\tabcolsep}{2.4pt}
\begin{tabular}{@{}lccccccc@{}}
\toprule
Method & \multicolumn{3}{c}{WebQA acc$_{\mathrm{approx}}$} & \multicolumn{3}{c}{OK-VQA} & Gap \\
\cmidrule(lr){2-4}\cmidrule(lr){5-7}
 & Score & vs.\ strong & vs.\ static & Score & vs.\ strong & vs.\ static & \\
\midrule
\texttt{direct} & \ph{} & \ph{} [\ph{},\,\ph{}] & \ph{} [\ph{},\,\ph{}] & \ph{} & \ph{} [\ph{},\,\ph{}] & \ph{} [\ph{},\,\ph{}] & \ph{} \\
\texttt{static\_text} & \ph{} & \ph{} [\ph{},\,\ph{}] & \ph{} [\ph{},\,\ph{}] & \ph{} & \ph{} [\ph{},\,\ph{}] & \ph{} [\ph{},\,\ph{}] & \ph{} \\
\texttt{static\_image} & \ph{} & \ph{} [\ph{},\,\ph{}] & \ph{} [\ph{},\,\ph{}] & \ph{} & \ph{} [\ph{},\,\ph{}] & \ph{} [\ph{},\,\ph{}] & \ph{} \\
\texttt{static\_mixed} & \ph{} & \ph{} [\ph{},\,\ph{}] & \ph{} [\ph{},\,\ph{}] & \ph{} & \ph{} [\ph{},\,\ph{}] & \ph{} [\ph{},\,\ph{}] & \ph{} \\
\texttt{rerank\_static} & \ph{} & \ph{} [\ph{},\,\ph{}] & \ph{} [\ph{},\,\ph{}] & \ph{} & \ph{} [\ph{},\,\ph{}] & \ph{} [\ph{},\,\ph{}] & \ph{} \\
NTE trigger & \ph{} & \ph{} [\ph{},\,\ph{}] & \ph{} [\ph{},\,\ph{}] & \ph{} & \ph{} [\ph{},\,\ph{}] & \ph{} [\ph{},\,\ph{}] & \ph{} \\
SeaKR & \ph{} & \ph{} [\ph{},\,\ph{}] & \ph{} [\ph{},\,\ph{}] & \ph{} & \ph{} [\ph{},\,\ph{}] & \ph{} [\ph{},\,\ph{}] & \ph{} \\
v1 $\Uomni$ gate & \ph{} & \ph{} [\ph{},\,\ph{}] & \ph{} [\ph{},\,\ph{}] & \ph{} & \ph{} [\ph{},\,\ph{}] & \ph{} [\ph{},\,\ph{}] & \ph{} \\
UniversalRAG & \ph{} & \ph{} [\ph{},\,\ph{}] & \ph{} [\ph{},\,\ph{}] & \ph{} & \ph{} [\ph{},\,\ph{}] & \ph{} [\ph{},\,\ph{}] & \ph{} \\
v2-Rule & \ph{} & \ph{} [\ph{},\,\ph{}] & \ph{} [\ph{},\,\ph{}] & \ph{} & \ph{} [\ph{},\,\ph{}] & \ph{} [\ph{},\,\ph{}] & \ph{} \\
v2-Cal & \ph{} & \ph{} [\ph{},\,\ph{}] & \ph{} [\ph{},\,\ph{}] & \ph{} & \ph{} [\ph{},\,\ph{}] & \ph{} [\ph{},\,\ph{}] & \ph{} \\
Oracle & \ph{} & \ph{} [\ph{},\,\ph{}] & \ph{} [\ph{},\,\ph{}] & \ph{} & \ph{} [\ph{},\,\ph{}] & \ph{} [\ph{},\,\ph{}] & \ph{} \\
\bottomrule
\end{tabular}
\end{table*}

\subsection{Ablations}
\label{sec:ablation}
Table~\ref{tab:ablation} removes one v2-Cal feature at a time, compares v2-Rule with v2-Cal, and leaves a row for the fixed $\delta=0.10$ and $\epsilon=1.0$.
The $\delta$ row is a development-split description.
It is not a point selected on the test set.
Refitting uses the calibration split only.
Every entry is \ph{}.

\begin{table}[t]
\caption{Ablations on the frozen test metrics after calibration-split refits.
WebQA is acc$_{\mathrm{approx}}$.
OK-VQA is official VQA accuracy.
Utility is the mean paired gap to the strong baseline on the calibration split (Section~\ref{sec:risk}), and is not an accuracy.
Every entry is \ph{}.
The last row is not used to pick the test operating point.}
\label{tab:ablation}
\scriptsize
\setlength{\tabcolsep}{3pt}
\begin{tabular}{@{}lccc@{}}
\toprule
Variant & WebQA & OK-VQA & Utility \\
\midrule
v2-Cal, six features & \ph{} & \ph{} & \ph{} \\
Without direct margin & \ph{} & \ph{} & \ph{} \\
Without $a^{\star}$ margin & \ph{} & \ph{} & \ph{} \\
Without margin gap & \ph{} & \ph{} & \ph{} \\
Without token F1 to $a^{\star}$ & \ph{} & \ph{} & \ph{} \\
Without KL vs.\ direct & \ph{} & \ph{} & \ph{} \\
Without retriever top-1 & \ph{} & \ph{} & \ph{} \\
v2-Rule & \ph{} & \ph{} & \ph{} \\
Development $\delta=0.10$, $\epsilon=1.0$ & \ph{} & \ph{} & \ph{} \\
\bottomrule
\end{tabular}
\end{table}

\subsection{Open Retrieval, Generalization, and Cost}
\label{sec:extra}
Table~\ref{tab:open} is the open-retrieval image setting.
The current interface builds it with stage~s2, not with a frozen stage-s4 test, so the table is not a test result.
Generalization to MMQA and to video is \ph{} and has no command in that interface.
Table~\ref{tab:cost} reserves mean generator calls, input tokens, and latency.
Planning estimates of GPU time are not entered here.

\begin{table}[t]
\caption{Open-retrieval WebQA image setting (\texttt{webqa\_open\_img}), stage~s2 interface only.
Not a frozen test.
Primary metric acc$_{\mathrm{approx}}$, plus normalized BART and QA$_{\mathrm{approx}}$.
Every entry is \ph{}.}
\label{tab:open}
\small
\begin{tabular}{@{}lccc@{}}
\toprule
Action & acc$_{\mathrm{approx}}$ & norm.\ BART & QA$_{\mathrm{approx}}$ \\
\midrule
\texttt{direct} & \ph{} & \ph{} & \ph{} \\
\texttt{static\_text} & \ph{} & \ph{} & \ph{} \\
\texttt{static\_image} & \ph{} & \ph{} & \ph{} \\
\texttt{static\_mixed} & \ph{} & \ph{} & \ph{} \\
\texttt{rerank\_static} & \ph{} & \ph{} & \ph{} \\
\bottomrule
\end{tabular}
\end{table}

\begin{table}[t]
\caption{Cost placeholders on the stage-s4 runs.
Calls, input tokens, and latency are \ph{}. The deployment row prefills only what the frozen rule reads and decodes the chosen action (Section~
ef{sec:cost}).}
\label{tab:cost}
\small
\begin{tabular}{@{}lccc@{}}
\toprule
System & Calls & Input tokens & Latency \\
\midrule
\texttt{direct} only & \ph{} & \ph{} & \ph{} \\
Always \texttt{static\_mixed} & \ph{} & \ph{} & \ph{} \\
SeaKR & \ph{} & \ph{} & \ph{} \\
v2, all actions decoded & \ph{} & \ph{} & \ph{} \\
v2-Rule, deployment path & \ph{} & \ph{} & \ph{} \\
\bottomrule
\end{tabular}
\end{table}

\subsection{Soft control: attention steering moves attention, not answers}
\label{sec:steer}
This subsection is the evidence for choosing hard control (Section~\ref{sec:softhard}).
Every number in this subsection is copied from the root-cause report dated 13~October~2026.
The measurements come from the legacy pipeline, which has known prompt and format problems; they motivate the analysis and are not results of the method.
They are not v2 results, and they are not taken from \texttt{offline\_v2}.

On a 64-question probe, the L1 distance of attention toward the supporting block moved from 0.77 to 0.26, and the answer flipped on 2 of 64 questions.
On the 38 questions in that probe where the gate was on, the median gap between the top-1 and top-2 first-token logits was about 3.4--3.5, and the median largest logit displacement caused by the bias was about 2.0.
On 29 of those 38 questions the gap exceeded the displacement.
The two flips had gaps of 0 and 0.75.
First-token predictive entropy was about 0.01--0.02, and on 65\% of questions $h_{\mathrm{sem}}<0.01$.
Before the bias, about 59\% of the attention already lay on the supporting evidence.
The bias raised that absolute mass by 0.011 on average.
The largest bias was about 1.2 nats (median of the per-question maximum 0.78, 90th percentile 1.07).
Among 306 gated-on questions, the median attention odds between the best and worst chunk were 1.185, against the propositional upper bound $e^{1.2}=3.32$ at $U=1$ and a full relevance gap.
The same gap is why a uniform span bias in the style of PASTA, or a positional recalibration in the style of Found-in-the-Middle, has little room to change an answer the first token has already determined.
This draft does not report a new PASTA or Found-in-the-Middle run.

On 512 questions decoded under one protocol on CUDA, the first token changed on 14 questions.
EAB-P versus static retrieval changed official QA by +1.30\%, with a 95\% interval of $[-0.11,+0.87]$ points, which includes zero.
The same comparison on MPS was +1.20\% and also includes zero.
Of a +0.304 point QA difference, two questions contributed +0.195 and +0.122, and the other 510 questions summed to $-0.013$.
Replacing CUDA with MPS and adding no bias changed 62 of 512 answers.
EAB-P changed 55--64 answers.

The bias formula applies on rows $a\ge\tend$ (implementation name \texttt{ans1}).
A previous default applied the bias only on decode rows (\texttt{ans2}) and was labeled as the paper form.
That label is withdrawn.
The gate in those runs read $U_{\mathrm{post}}$, not the retrieval-time uncertainty named in the text.
Appendix~\ref{app:dca} records the local algebra of this retired bias.
It is not an accuracy result.

\subsection{Failure Cases and Format Sensitivity}
\label{sec:cases}
Failure counts for missing evidence, wrong modality, and format mismatch are \ph{}.
No case is filled in, because the v2 runs have not been executed.
Format sensitivity uses the single frozen prompt from Section~\ref{sec:setup}.
A comparison across prompt variants, if it is run, is a development-split description and is \ph{} here.
It is not a search over prompts on the test split.

%% ===============================================================
\section{Conclusion}
We describe \method{} v2, a per-question choice among direct answering, static text, static image, static mixed evidence, and static evidence reranked by a cross-encoder.
v2-Rule fits no weights.
v2-Cal is one small ridge regression per evidence action.
Fixed-sequence one-sided paired tests at $\delta=0.10$ and $\epsilon=1.0$ supply the expected-score guarantee in Section~\ref{sec:risk}.
The primary comparison uses the strongest single action on the calibration split as the denominator.
The measurement, the main table, and the ablations have not been run, and their numbers are \ph{}.
The attention-steering subsection quotes the 13~October~2026 root-cause report: attention moves, and answers do not.
Appendix~\ref{app:limitations} states the limitations.


\bibliographystyle{ACM-Reference-Format}
\bibliography{references}

\appendix
%% ===============================================================
\section{Evaluation Details}
\label{app:protocol}
This appendix restates the v2 protocol.
The v2 runs have not been executed.
Placeholder cells in the body stay \ph{}.

\textit{Splits.}
WebQA draws at least 500 calibration questions and 256 development questions from the training split.
Those two sets are disjoint from each other and from the test split.
The test split is all 4{,}966 validation questions.
Calibration is not taken from the WebQA validation split.
OK-VQA uses 500 calibration questions, 256 development questions, and the remainder of the validation split.
The three OK-VQA sets are pairwise disjoint.
An empty test set is an error.
For a validation split of 5{,}046 questions the remainder is 4{,}290.
Question-id lists are hashed, and the test command checks the hash of the frozen decision file \texttt{runs/v2\_decisions/}.

\textit{Actions.}
The main-method actions are \texttt{direct}, \texttt{static\_text}, \texttt{static\_image}, \texttt{static\_mixed}, and \texttt{rerank\_static}.
They share one prompt, one decoding budget, and one candidate pool inside a run.
The prompt asks for the frozen answer format.
It does not ask for a sources line in place of an answer.
WebQA uses one sentence aligned with the reference style.
OK-VQA uses at most three words and official VQA normalization.
The prompt is fixed before any model run and is the same for every arm.
Evidence Attention Bias is not one of the five actions.

\textit{Metrics and baselines.}
OK-VQA reports official VQA accuracy only.
WebQA's primary metric is acc$_{\mathrm{approx}}$ (Appendix~\ref{app:webqa_scoring}), with normalized BART and QA$_{\mathrm{approx}}$ on the full set and on image and text subsets.
The primary WebQA setting is the distractor pool.
Open retrieval is the second setting.
The strong baseline is the single action with the best calibration-set score.
Primary relative gain uses that score as the denominator.
Gain versus \texttt{static\_mixed} is also reported.
The oracle is the per-question best action.
Adaptive comparators (NTE trigger, SeaKR, the v1 $\Uomni$ gate, UniversalRAG) are not candidates for the strong baseline.
\texttt{gold-in-top-$k$} is offline analysis only.

\textit{Statistics and selection.}
Intervals are paired bootstrap 95\% confidence intervals over questions, with Holm correction for multiple comparisons.
Greedy decoding does not contribute a seed standard deviation.
The risk-control claim uses the mean paired gap to the strong baseline (Section~\ref{sec:risk}).
Thresholds and grids are chosen on the calibration or development split.
The test split is scored once.
A grid is not selected on the test split.
The development gates in Section~\ref{sec:setup} are criteria.
This draft does not claim that a gate has passed.

%% ===============================================================
\section{Local Algebra of the Retired Attention Bias}
\label{app:dca}
The identities below describe the attention bias removed from the main method.
They are not v2 accuracy results, and they do not say that the bias improves answers or order robustness.
Section~\ref{sec:steer} quotes the empirical measurement.

Fix one layer, one head, and one attention row, with pre-bias attention mass $\mu_j$ on chunk $j$.
Retrieved chunks $j\in\{1,\dots,k\}$ receive a logit bias $\lambda r_j$ on the rows $a\ge\tend$ (\texttt{ans1}).
Group $0$ is the remaining visible keys, with $r_0=0$.
When the old gate is active, $\lambda=\gamma\Uomni$; otherwise $\lambda=0$.
The previous default path that biased only decode rows (\texttt{ans2}) is not this formula, and it is not labeled as the paper form.
The gate that produced the quoted runs read $U_{\mathrm{post}}$.

\begin{proposition}[Bounded relevance tilt]
\label{thm:tilt}
Under this bias, attention inside each span keeps its proportions, and the chunk masses become
\begin{equation}
\mu_j'=\frac{\mu_j e^{\lambda r_j}}{\sum_{l=0}^{k}\mu_l e^{\lambda r_l}}.
\label{eq:local}
\end{equation}
For any two chunks, $\mu'_j/\mu'_l=(\mu_j/\mu_l)\,e^{\lambda(r_j-r_l)}$ with $|\lambda(r_j-r_l)|\le\gamma$.
\end{proposition}

\begin{proof}
A key $i$ in $\mathrm{span}(d_j)$ receives logit $s_i+\lambda r_j$, so its weight is $\alpha_i e^{\lambda r_j}/\sum_{i'}\alpha_{i'}e^{\lambda r(i')}$.
Keys in one span share $e^{\lambda r_j}$, so proportions inside the span are unchanged.
Summing over the span gives Equation~\eqref{eq:local} with $Z(\lambda)=\sum_l\mu_l e^{\lambda r_l}$.
For any distribution $p$ over the groups,
\[
\lambda\sum_j p_j r_j-D_{\mathrm{KL}}(p\,\|\,\mu)
=\log Z(\lambda)-D_{\mathrm{KL}}(p\,\|\,\mu'),
\]
which is maximized uniquely at $p=\mu'$.
The odds ratio is the quotient of two instances of Equation~\eqref{eq:local}.
$|\lambda(r_j-r_l)|\le\gamma$ because $\lambda\le\gamma\Uomni\le\gamma$ and $r\in[0,1]$.
\end{proof}

Chunk replication with $n_j=\max(1,\mathrm{round}(e^{\gamma\Uomni r_j}))$ copies matches the same reweighting at bias $\log n_j$, up to rounding, under three local assumptions: replica logits at the row are unchanged, values are unchanged, and every copy is visible to the row.
Rotary embeddings, causal masking, and cross-layer residuals can break those assumptions.
A full-system accuracy comparison of replication and the bias is not a v2 result and is not reported.

\begin{proposition}[Local copy identity]
\label{thm:dca-local}
Under the three assumptions, pooling the copies of key $i\in\mathrm{span}(d_j)$ gives the EAB weight at bias $b_j=\log n_j$, and the chunk masses match.
\end{proposition}

\begin{proof}
Both normalizers equal
\begin{equation}
D=\sum_{i\in\text{group }0}e^{s_i}+\sum_{\ell=1}^{k}n_\ell\sum_{i\in\mathrm{span}(d_\ell)}e^{s_i},
\label{eq:dca-norm}
\end{equation}
because $n_j e^{s_i}=e^{s_i+\log n_j}$.
Pooling copies gives weight $n_j e^{s_i}/D$.
Equal pooled weights and equal values give the same attention output at the row.
\end{proof}

%% ===============================================================
\section{Supplementary Placeholders}
\label{app:supp}
Open-retrieval scores, cost, and the paired-gap intervals are the placeholder tables in Section~\ref{sec:extra} and Table~\ref{tab:main}.
No seed-wise standard deviation is attached to a greedy run.
Sensitivity to $\delta$ is the development-split row of Table~\ref{tab:ablation}.
That row is not a test-set search.
Backbone transfer, if it is run, is \ph{} and is not claimed here.
Secondary WebQA numbers for a filled run are normalized BART and QA$_{\mathrm{approx}}$ from the same predictions as acc$_{\mathrm{approx}}$.
They are \ph{} until stage~s4 is run.

%% ===============================================================
\section{WebQA Scoring}
\label{app:webqa_scoring}
The published WebQA evaluator~\cite{chang2022webqa} computes acc$_{\mathrm{approx}}$ from category-specific F1 for color, shape, number, and yes/no questions and from recall for the other categories.
We combine this accuracy component $a_i$ with the ParaBank2 BARTScore checkpoint of the WebQA authors.
For reference $r_i$, answer $y_i$, and reference-to-answer log score $s$,
\begin{equation}
F_i=\min\bigl\{e^{s(r_i,y_i)-s(r_i,r_i)},1\bigr\},\qquad
\mathrm{QA}_{\mathrm{approx}}=\frac{100}{N}\sum_{i=1}^{N}a_i F_i .
\label{eq:webqa}
\end{equation}
$F_i$ is the normalized BART factor.
Public validation has reference answers and does not have the keyword labels of the hidden test set, which is why the combined score is called QA$_{\mathrm{approx}}$.
The primary reported WebQA number is acc$_{\mathrm{approx}}$.
Normalized BART and QA$_{\mathrm{approx}}$ are reported with it, including image and text subsets.
Loose lexical scores are not reported beside them.
BART scoring uses CPU FP32, eager attention, and one checkpoint for every action.
Filled values of these metrics are \ph{} until the v2 run exists.
The 512-question static figures in Section~\ref{sec:setup} are quotations from the root-cause report, not cells of Table~\ref{tab:webqa}.

%% ===============================================================
\section{Error and Case Analysis}
\label{app:error}
Shares of errors from missing top-$k$ evidence, from the wrong modality, and from answer normalization are \ph{}.
The denominator, once a run exists, is the number of scored test questions, and failures to decode are counted in that denominator.
No example is inserted in this draft.
The attention-bias cases discussed in Section~\ref{sec:steer} are the root-cause report's cases, not a new error study.

%% ===============================================================
\section{Limitations and Ethical Considerations}
\label{app:limitations}
\textit{Limitations.}
v2 decodes several actions per question.
Thresholds and the ridge coefficients of v2-Cal depend on the calibration split.
The strongest single action is chosen on that same sample, which makes the risk claim slightly anti-conservative.
The guarantee in Section~\ref{sec:risk} bounds the expected score of the chosen policy, under i.i.d. questions and an asymptotically valid test, at $\delta=0.10$ and a tolerance of $\epsilon=1.0$ point.
The image-evidence pipeline is a separate 50-question diagnostic whose outcome is \ph{}.
\texttt{gold-in-top-$k$} is not available as an inference-time feature.
MMQA and video are outside the current command interface.
Historical cached baselines, loose WebQA scores, greedy seed standard deviations, and test-set grid selection are withdrawn and are not repaired by filling \ph{}.

\label{app:ethics}
\textit{Ethical considerations.}
Retrieved Web text and images can be wrong, biased, or adversarial, and a selector that sometimes trusts them can repeat that content.
The direct action is kept in the action set so that retrieval is not mandatory.
Datasets used here are public.
Deployments on private content need access controls.
\tbd{[[AI-use disclosure]]}


\end{document}
